% Replacement algorithm and FPGA architecture sections.
% Required packages: amsmath, amssymb, graphicx, booktabs, tikz.
% Required TikZ libraries: arrows.meta, positioning, calc.
% Compile from the manuscript directory so figures/ resolves correctly.
% This file makes no claim of measured physical stabilization or throughput.

\section{Gravity-Referenced Camera Correction}
\label{sec:gl_algorithm}
\label{sec:algorithm}

\subsection{Scope and Calibration Contract}
The method separates static optical mapping, time-varying camera orientation,
and frame association. An archived inverse map accounts for the fixed lens,
cylindrical projection, crop, and strip placement. A low-rate orientation
engine converts one rig roll/pitch pair into six camera-specific corrections.
The renderer applies these corrections to the map segments associated with
one leased source frameset. The FPGA receives orientation estimates; it does
not estimate attitude from raw inertial samples. There is no live yaw input,
heading lock, translation estimate, or rolling-shutter model in this path.

For pure rotation and a calibrated central camera, a rotation in ray space
can be mapped back to a cylinder. Thus a sufficiently general global
spherical or cylindrical warp can compensate rotation. The practical issue
here is narrower: a common vertical shift or planar roll applied to an
already assembled panorama does not generally reproduce the different local
corrections around the ring. The proposed implementation uses camera-specific
affine corrections to retain the existing segment-based mapping and sampling
hardware. It is an approximation to ray reprojection, rather than a claim
that all global rotational correction is impossible.

Let camera $i$ have mounting yaw $\psi_i$, pitch $p_i^0$, and roll $r_i^0$.
The six EO yaws are $(0,60,120,180,240,300)^\circ$; the IR yaws are
$(90,150,210,270,330,390)^\circ$. The $390^\circ$ value is retained in
calibration and normalized modulo $360^\circ$ for trigonometric lookup.
Initial mounting roll and pitch are zero for all cameras. Subsequently,
per-camera UI roll/pitch trims are retained across boots. The MCU applies
the selected mounting profile once, recording revision \texttt{0x7106B1F1}
in a CRC-protected two-page journal. Runtime yaw edits are stored by the MCU,
but the present FPGA uses the fixed yaw constants above. Only roll/pitch
trims feed its orientation engine.

Table~\ref{tab:gl_calibration} distinguishes the installed optical profile
from the selected mounting-profile INI. In particular, adopting the mounting
angles does not regenerate the optical maps. The INI's example live poses
$(0.75,-5.00)^\circ$ for EO and $(1.25,6.50)^\circ$ for IR are test inputs,
not runtime defaults. Likewise its output-mode and DDR-offset fields are
not the hardware mode/address contract.

\begin{table*}[t]
\caption{Calibration and preserved optical configuration. The ``selected INI''
column is recorded for reproducibility; only its mounting arrays are imported.
Focal lengths and principal points are in original source pixels.}
\label{tab:gl_calibration}
\centering\footnotesize\setlength{\tabcolsep}{4.3pt}
\begin{tabular}{@{}lccc@{}}
\toprule
Parameter & Installed EO & Installed IR & Selected INI: EO / IR \\
\midrule
Source size; cameras & $1920\!\times\!1080$; 6 & $640\!\times\!512$; 6 & unchanged \\
$(f_x,f_y)$ & $(1141.834740,1142.434960)$ & $(564.835060,564.024805)$ & unchanged \\
Normalized principal point & $(0.489836,0.514300)$ & $(0.497330,0.476162)$ & unchanged \\
Distortion principal point & $(940.486,555.444)$ & $(318.291416,243.794865)$ & unchanged \\
$(k_1,k_2,k_3)$ & $(-.163318,.165940,-.064582)$ & $(-.327245098,.299803510,-.197210469)$ & unchanged \\
$(p_1,p_2)$ & $(.000236,.000773)$ & $(.000324600,.000015975)$ & unchanged \\
Map resize factor & $0.5$ & $1.0$ & $0.56296$ / $1.0$ \\
Cylindrical horizontal/vertical FOV & $57.56^\circ/48.56^\circ$ & $54.5^\circ/45.0^\circ$ & $77^\circ/54^\circ$ / $65^\circ/52^\circ$ \\
Crop height/width scale & $0.888/1.0$ & $0.9375/1.0$ & $0.79/0.70$ / $0.65/0.9$ \\
Map-profile overlap; feather strength & $50$; $1.0$ & $30$; $1.0$ & $60$; $1.0$ / $30$; $1.0$ \\
Logical panorama size & $3840\!\times\!480$ & $3576\!\times\!480$ & $3840\!\times\!480$ / $3240\!\times\!480$ \\
Source pitch shift, pixels/degree & $-9.80$ & $-8.72$ & newer C: $-17.408$ / $-8.72$ \\
Per-camera strip widths & $650,655,655,655,655,655$ & $616,621,640,621,621,621$ & installed segment ROMs retained \\
Panorama strip starts & $0,633,1271,1909,2547,3185$ & $0,587,1179,1771,2363,2955$ & installed placement retained \\
\bottomrule
\end{tabular}
\end{table*}

\subsection{Upstream Attitude Estimate}
The copied MCU firmware parses IIM-46234 accelerometer, gyroscope, and
temperature packets. Its requested sensor profile is a $1000$ Hz packet
rate (divider $1$), output selection \texttt{0x07}, and bandwidth setting
\texttt{0x44} (BW4 for both sensors). These are configuration values, not
measured acquisition rates or a measured filter response. Eight accepted
samples are averaged before each six-axis Madgwick update~\cite{madgwick2011gradient}, giving a nominal
$125$ Hz fusion rate. The gain is $\beta=0.12$, the initial quaternion is
$(1,0,0,0)$, and the configured gyro-bias vector is
$(0.421587,-0.133538,0.064359)$ degrees/s. Gyro values are converted to
radians/s after bias subtraction; acceleration and gyro axes are passed
directly to fusion, without a calibrated rig-axis remapping in this path.
The update uses gyro quaternion integration, normalized accelerometer
gradient feedback, and quaternion normalization; no magnetometer is used.
An isolated host test of the copied filter at identity orientation, zero
gyro, and acceleration $(0,0,1)$ yields a non-finite quaternion because
zero-gradient normalization is unguarded; this is a numerical boundary
finding, not an observed physical failure.
The assumed packet rate is refreshed from groups of $32$ timestamp/counter
intervals, and fusion uses that rate divided by eight.

Quaternion-derived roll/pitch are clipped to $\pm90^\circ$, smoothed, then
rounded to centidegrees and quantized in $10$ cd steps. The implemented
output smoothing is $z\leftarrow z+\alpha(\theta-z)$, with
\begin{equation}
 \alpha=\operatorname{clip}\!\left(
 \frac{1000/f_s}{10+1000/f_s},0.02,1\right),
\end{equation}
where $f_s$ is the estimated packet rate in Hz and times are in ms.
The shadow angles update every ten accepted packets or when a VSYNC request
is pending. Since this schedule differs from the interval used to calculate
$\alpha$, the constant $10$ is not a measured end-to-end time constant.
PC4 VSYNC sets a pending flag; the main loop services it for panorama modes
\texttt{0x14}/\texttt{0x15} and sends the current shadow pair to the FPGA.
The source packet timestamp is not included in that transaction. A separate
$250$ ms no-DRDY status timeout also does not establish exposure alignment.

\subsection{Azimuth-Dependent Orientation Projection}
All angle interfaces use integer centidegrees (cd), with $100$ cd per degree.
Using the usual right-handed axis-rotation matrices, the implemented order is
\begin{align}
 R_s &= R_x(r)R_y(p), & T_i &= R_z(\psi_i)R_y(p_i^0),\\
 R_{m,i} &= T_iR_x(r_i^0), & R_i &= R_sR_{m,i}.
 \label{eq:gl_matrix_chain}
\end{align}
These are four distinct fixed-point matrix products. Their order and
intermediate quantization are part of the numerical specification.
For zero mounting trims and small $r,p$ in radians,
\begin{equation}
 \begin{bmatrix}r_i\\p_i\end{bmatrix}
 \simeq
 \begin{bmatrix}\cos\psi_i&\sin\psi_i\\-\sin\psi_i&\cos\psi_i\end{bmatrix}
 \begin{bmatrix}r\\p\end{bmatrix}.
 \label{eq:gl_small_angle}
\end{equation}
This expression explains the azimuth dependence in
Fig.~\ref{fig:gl_projection}; the hardware evaluates
\eqref{eq:gl_matrix_chain}, not this linearization. Axis names and signs here
define the software convention. Their correspondence to the physical IMU
and rig still requires installation-specific measurement.

\begin{figure}[t]
\centering% Native vector artwork; requires tikz with arrows.meta and calc.
\begin{tikzpicture}[x=1cm,y=1cm,>=Latex,font=\footnotesize,
  glgbox/.style={draw,rounded corners=1pt,align=center,inner sep=5pt}]
  \draw[gray] (1.45,0) circle (1.08);
  \draw[->] (1.45,0)--(2.8,0) node[right] {$x$};
  \draw[->] (1.45,0)--(1.45,1.35) node[above] {$y$};
  \fill (1.45,0) circle (1.4pt);
  \foreach \a/\n in {0/0,60/1,120/2,180/3,240/4,300/5}{
    \draw[->,thick] (1.45,0)--++(\a:1.08);
    \node[fill=white,inner sep=1pt] at ($(1.45,0)+(\a:1.28)$) {\n};
  }
  \draw[->] (2.06,0) arc[start angle=0,end angle=60,radius=.61];
  \node at (2.18,.40) {$\psi_i$};
  \node[align=center] at (1.45,-1.80) {EO azimuths shown\\IR ring: $+90^\circ$};
  \node[glgbox,text width=4.15cm] (mount) at (5.60,.82)
    {Mounting rotation\\$R_{m,i}=R_z(\psi_i)R_y(p_i^0)R_x(r_i^0)$};
  \node[glgbox,text width=4.15cm] (body) at (5.60,-.40)
    {Measured rig tilt\\$R_i=[R_x(r)R_y(p)]R_{m,i}$};
  \draw[->] (mount)--(body);
  \node[glgbox,text width=4.15cm] (local) at (5.60,-1.73)
    {Local Euler correction\\$r_i\simeq r\cos\psi_i+p\sin\psi_i$\\
     $p_i\simeq p\cos\psi_i-r\sin\psi_i$};
  \draw[->] (body)--(local);
  \node[align=center,text width=7.8cm] at (4.0,-2.85)
    {Small-angle expressions explain azimuth dependence.\\
     The implementation evaluates the quantized matrix products.};
\end{tikzpicture}

\caption{Camera mounting orientation changes the local roll/pitch induced by
one rig tilt. The small-angle relation is explanatory; finite-angle values
come from the fixed-point matrix chain.}
\label{fig:gl_projection}
\end{figure}

The orientation trigonometric ROM contains $36000$ sine/cosine pairs at
$0.01^\circ$ spacing. Its scale is $32767$, including the matrix identity
entries. Each dot product uses the quantizer
\begin{equation}
 \begin{aligned}
 Q(a)&=\operatorname{sat}_{[-32768,32767]}
       \left(\left\lfloor\frac{a+2^{14}}{2^{15}}\right\rfloor\right),
 \\
 [AB]_{jk}&=Q\!\left(\sum_{\ell=0}^2 A_{j\ell}B_{\ell k}\right).
 \end{aligned}
 \label{eq:gl_mat_quant}
\end{equation}
The positive bias applies even to negative accumulators. Replacing it with
symmetric rounding, floating-point multiplication, or a reassociated matrix
chain changes the bit-level result. Products are signed 32-bit values and
the RTL accumulator is signed 36-bit. ROM generation uses the C library's
nearest rounding with half-way cases away from zero.

The inverse trigonometric operator $A(y,x)$ implements integer
$\operatorname{atan2}$ in centidegrees. For $m=\min(|x|,|y|)$ and
$d=\max(|x|,|y|)>0$, it indexes a $4097$-entry table by
\begin{equation}
 \begin{aligned}
 j&=\left\lfloor\frac{4096m+\lfloor d/2\rfloor}{d}\right\rfloor,\\
 L[j]&=\operatorname{round}_{\mathrm{away}}
       \left(\frac{18000}{\pi}\tan^{-1}\frac{j}{4096}\right).
 \end{aligned}
\end{equation}
Axis exchange and quadrant reconstruction give the complete angle;
$A(0,0)=0$. With $v=\operatorname{clip}(-R_{i,20},-32767,32767)$,
\begin{equation}
 p_i=A\!\left(v,\left\lfloor\sqrt{2^{30}-v^2}\right\rfloor\right).
\end{equation}
If the magnitude of the ROM cosine of $p_i$ exceeds $8$, then
$r_i=A(R_{i,21},R_{i,22})$ and
$y_i=A(R_{i,10},R_{i,00})$; otherwise $r_i=0$ and
$y_i=A(-R_{i,01},R_{i,11})$. The extracted yaw is not used by the renderer.
The $32767$ scale and integer square root mean that ideal right-angle
inputs need not yield exactly $9000$ cd.

\subsection{Correction Range and Preserved Pixel Mapping}
The controller first computes all six camera roll/pitch pairs and
$M=\max_i\{|r_i|,|p_i|\}$. When $(r,p)\ne(0,0)$ and $M>1000$, it replaces
each live component $v\in\{r,p\}$ by
\begin{equation}
 v'=\operatorname{sgn}(v)
   \left\lfloor\frac{1000|v|+\lfloor M/2\rfloor}{M}\right\rfloor
 \label{eq:gl_global_scale}
\end{equation}
and recomputes the six cameras once. Static trims are never rescaled.
This single global adjustment is distinct from independently clipping six
camera poses, and it need not force every recomputed angle within the limit.
A separate flag reports that condition. The preserved renderer finally
clips each camera roll and pitch to $[-1000,1000]$ cd. Disabled or unusable
live correction bypasses the projection and returns the static trim arrays
exactly; enabled zero tilt still executes the quantized matrix chain.

Let $\boldsymbol b(u,v)$ be the archived inverse map, in original source
pixels, with source center $(c_x,c_y)=(W/2,H/2)$. For the final camera
angles $\bar r_i,\bar p_i$ in degrees, the source coordinate is
\begin{align}
 t_i&=\operatorname{round}_{\mathrm{away}}(K_p\bar p_i/2^{16}),\\
 d_x&=b_x-c_x, & d_y&=b_y+t_i-c_y,\\
 x_s&=c\,d_x+s\,d_y+c_x,&
 y_s&=-s\,d_x+c\,d_y+c_y,
 \label{eq:gl_map_affine}
\end{align}
where $c=\cos(-\bar r_i)$, $s=\sin(-\bar r_i)$ and
$K_p=\operatorname{round}_{\mathrm{away}}(2^{16}k_p)$ is $-642253$ for EO
($k_p=-9.80$) or $-571474$ for IR ($k_p=-8.72$).
Translation precedes rotation.
This pixel path intentionally uses a different trigonometric scale,
$32768$, so zero angle is an exact identity. Its $1001$-entry magnitude
ROM covers $0$--$1000$ cd, restores signs at runtime, and stores the pitch
translation as a signed 8-bit integer ($\pm98$ EO or $\pm87$ IR rows).
The newer C optical path's resize-normalized EO factor is Q16
$-1140855$, approximately $-17.408$ pixels/degree; it is not substituted
into the retained renderer. Therefore exact agreement of the orientation
engine with C does not imply end-to-end pixel agreement with that newer
optical pipeline.

Map origins and running coordinates are signed 32-bit integers with
$16$ fractional bits. Each row segment has at most $64$ pixels; stored
16-bit increments have $5$ fractional bits for EO and $4$ for IR, and
are shifted left by $11$ or $12$ before transformation. Transformed
increments use signed 24-bit Q16 storage. Coordinate products are shifted
arithmetically by $15$ without an added rounding bias. Since the map is
affine within a segment, transformed endpoints bound every coordinate in it.

The sampler accepts $0\le x_s\le W-2$ and $0\le y_s\le H-2$; other
samples become black. It uses two vertically adjacent samples, rather than
general bilinear interpolation. EO quantizes the horizontal coordinate to
the output pixel's chroma parity and vertically interpolates luma and chroma
separately with a $2^{15}$ rounding bias before shifting by $16$. IR uses
$\lfloor x_s\rfloor$, truncating fixed-point vertical interpolation, and
constant chroma $128$. Archived seam-weight ROMs and strip placement remain
unchanged. EO overlaps contain $17$ output pixels; ordinary IR overlaps
contain $29$. The IR camera-2/3 fold seam uses the inherited $48$-pixel tail,
$90$-pixel head, and specialized alpha-index mapping.

\section{Frame-Coherent FPGA Architecture}
\label{sec:gl_architecture}
\label{sec:fpga}

\subsection{Atomic Pose Reception and Source-Epoch Association}
The runtime flow is shown in Fig.~\ref{fig:gl_runtime}.
The MCU writes I\textsuperscript{2}C slave address \texttt{0x36}, starting
at register \texttt{0x10}, with four little-endian bytes: signed 16-bit
roll followed by signed 16-bit pitch. The receiver commits only at STOP
after exactly four acknowledged sequential data bytes. Partial, excess,
misordered, or restarted bursts cannot publish a torn pose. A separate
one-byte transaction at \texttt{0x04} accepts \texttt{0xA1} to enable
correction and \texttt{0xA0} to disable it; reset is disabled. The MCU
enables correction only for UI On together with Attitude Sensor.
Request/acknowledge CDC mailboxes transfer coherent pose and control into
the DDR user-interface (UI) clock domain. The 32-bit receipt sequence counts
accepted FPGA pose transactions; it is not an IMU timestamp.

\begin{figure}[t]
\centering% Native vector artwork; requires tikz with arrows.meta and positioning.
\begin{tikzpicture}[>=Latex,font=\footnotesize,node distance=5mm,
  glfbox/.style={draw,rounded corners=1pt,align=center,text width=7.25cm,inner sep=4pt},
  glflink/.style={->,thick}]
  \node[glfbox] (receive) {Commit complete roll/pitch transaction at STOP\\
    Transfer coherent pose and local receipt sequence to UI clock};
  \node[glfbox,below=of receive] (capture) {On each family's trigger: capture pose, age, enable, revision\\
    Independent EO and IR histories; key = (family, epoch)};
  \node[glfbox,below=of capture] (lease) {Obtain source-bank lease; freeze static calibration\\
    Look up that family's epoch, even if the view has changed};
  \node[glfbox,below=of lease] (policy) {Found, fresh, enabled, and same calibration revision?\\
    Yes: project six cameras; optionally rescale live tilt once\\
    No: use exact static roll/pitch values};
  \node[glfbox,below=of policy] (publish) {Publish one six-camera angle set atomically\\
    Accept render request and freeze its angles and identity};
  \node[glfbox,below=of publish] (render) {Reuse matching tables or rebuild from the frozen set\\
    Replay leased pixels; render, blend, and retire the pass};
  \foreach \a/\b in {receive/capture,capture/lease,lease/policy,policy/publish,publish/render}
    \draw[glflink] (\a)--(\b);
  \draw[->,dashed] (render.east)--++(.30,0)|-(capture.east);
\end{tikzpicture}

\caption{Pose selection follows the source-frame lease. New poses may arrive
throughout rendering but cannot modify an accepted pass. A lease or mode
change before acceptance cancels pending computation.}
\label{fig:gl_runtime}
\end{figure}

EO strobe and IR genlock drive independent 16-bit epoch counters. Each
family's six capture writers receive its Gray-coded epoch; a synchronized
trigger edge also captures the latest coherent pose in that family's
16-entry history. Each entry stores the complete epoch tag, roll/pitch,
receipt sequence, pose age, validity, enable state, and calibration revision.
Both histories continue capturing while the other family is displayed.
Consequently identical numeric EO and IR epochs cannot select each other's
pose. Full tags reject an entry overwritten by an index alias; requests
must not remain live across an entire 16-bit epoch wrap.

Pose age is zero on its receipt edge and increases in UI cycles. Freshness
is evaluated \emph{at trigger capture}, with the inclusive default bound
$23{,}340{,}000$ cycles ($100$ ms at nominal $233.4$ MHz).
Age is 32-bit saturating, with a sticky overflow flag that forces unusable
status. A pose update on the capture edge takes precedence; capture also
takes precedence over a simultaneous lookup of the same tag. Neither this
age bound nor an epoch establishes the physical exposure time. IMU filtering,
delivery latency, trigger-to-exposure delay, and their variation remain
separate calibration quantities.

The frameset manager provides source banks sharing an epoch among the
present cameras and holds the panorama lease until retirement. The pose
controller looks up the key (sensor family, leased epoch), freezes the
current static trims, and checks the captured calibration revision. Missing,
stale, disabled, or revision-mismatched data selects the exact static bypass.
The static register-bank path requires $2^{25}$ unchanged UI cycles before
adopting new values and clamps them to $\pm1000$ cd; the MCU accepts UI
trims within $\pm500$ cd. The two 16-bit bank-update counters form the
32-bit revision, so a change in either bank invalidates old calibration
association. This settling filter is distinct from the atomic live-pose
transaction protocol.

\subsection{Frozen Parameters and Source-Driven Rendering}
Figure~\ref{fig:gl_hardware} shows the implemented data and control paths.
One iterative projection engine evaluates six cameras sequentially and
publishes their $6\times16$-bit roll and pitch arrays atomically. At render
acceptance, the renderer snapshots those arrays, sensor selection, and
revision. Table construction reads the same frozen registers. Identical
angles and sensor selection permit table reuse; updates arriving during a
build or active pass cannot invalidate it. Replay is held until the new
source requirements are valid, avoiding a first-row fetch based on an old
table. The lease/pose controller cancels obsolete work before acceptance;
accepted coefficients remain fixed through completion or abort.

\begin{figure*}[t]
\centering% Native vector artwork; requires tikz with arrows.meta, positioning, calc.
\begin{tikzpicture}[x=1cm,y=1cm,>=Latex,font=\footnotesize,
  glhbox/.style={draw,rounded corners=1pt,align=center,inner sep=4pt,minimum height=.90cm},
  glhdata/.style={->,thick},glhctrl/.style={->,dashed}]
  \node[glhbox,text width=2.20cm] (imu) at (1.25,0)
    {I$^2$C pose/control\\atomic receiver + CDC};
  \node[glhbox,text width=2.45cm] (hist) at (4.25,0)
    {EO / IR epoch histories\\16 full tags per family};
  \node[glhbox,text width=2.55cm] (pose) at (7.40,0)
    {Lease lookup + projection\\one frozen angle set};
  \node[glhbox,text width=2.35cm] (table) at (10.55,0)
    {Segment transform + sort\\tables rebuilt on change};
  \node[glhbox,text width=2.35cm] (rom) at (13.60,0)
    {Static map segments\\trig and seam ROMs};
  \draw[glhctrl] (imu)--(hist);
  \draw[glhctrl] (hist)--(pose);
  \draw[glhctrl] (pose)--(table);
  \draw[glhctrl] (rom)--(table);
  \node[align=center] (cal) at (7.40,1.0) {Mount yaw + UI trims / revision};
  \draw[glhctrl] (cal)--(pose);
  \node[glhbox,text width=2.20cm] (camera) at (1.25,-2.05)
    {Six cameras, one family\\capture writers};
  \node[glhbox,text width=2.45cm] (ddr) at (4.25,-2.05)
    {DDR source banks\\four banks per camera};
  \node[glhbox,text width=2.55cm] (replay) at (7.40,-2.05)
    {Frameset lease + replay\\global row/column bounds};
  \node[glhbox,text width=2.35cm] (ring) at (10.55,-2.05)
    {Input line ring\\three rows per camera};
  \node[glhbox,text width=2.35cm] (walk) at (13.60,-2.05)
    {Walker + sampler\\source-row scheduling};
  \foreach \a/\b in {camera/ddr,ddr/replay,replay/ring,ring/walk}
    \draw[glhdata] (\a)--(\b);
  \draw[glhctrl] (camera.north)--++(0,.45)-|(hist.south);
  \node[fill=white,inner sep=1pt] at (2.75,-1.19) {trigger epoch};
  \draw[glhctrl] (replay)--(pose);
  \draw[glhctrl] (table.south)--++(0,-.35)-|(walk.north);
  \node[glhbox,text width=2.35cm] (window) at (13.60,-4.05)
    {Output RowWindow\\256 EO / 512 IR rows};
  \node[glhbox,text width=2.35cm] (emit) at (10.55,-4.05)
    {Row-ready emitter\\seam blend + fold};
  \node[glhbox,text width=2.55cm] (out) at (7.40,-4.05)
    {DDR output banks\\build / pending / scanout};
  \node[glhbox,text width=2.45cm] (scan) at (4.25,-4.05)
    {Raster output\\may repeat completed image};
  \foreach \a/\b in {walk/window,window/emit,emit/out,out/scan}
    \draw[glhdata] (\a)--(\b);
  \draw[glhctrl] (table.west)--++(-.35,0)--++(0,-3.30)-|(emit.north);
  \node[align=center] at (7.40,-5.05)
    {Solid: pixels / image storage\quad Dashed: pose / scheduling\quad EO and IR share the renderer};
\end{tikzpicture}

\caption{Implemented panorama path. Full source-frame address spaces reside
in DDR even when EO capture writes a proven subwindow. On-chip input rows
are replay caches, not the complete image storage. The same processing
engine serves EO or IR; output scanout is decoupled from image production.}
\label{fig:gl_hardware}
\end{figure*}

The base map contains $480\times6\times11=31680$ EO segments or
$480\times6\times10=28800$ IR segments. The table builder transforms each
segment, chooses traversal direction so source $y$ is nondecreasing, and
counting-sorts entries by their first source row. It stores a bucket end
index for each source row and a row-ready threshold equal to the maximum
last source row needed by each output row. A $4096$-entry queue retains
partially traversed segments. For source row $y_0$, the walker drains the
previously queued work and newly eligible segments, emits their pixels,
and requeues unfinished segments. Coordinates advance by addition.

DDR replay fills a three-row ring for each camera: rows $y_0$ and $y_0+1$
are sampled while $y_0+2$ is prefetched into the third slot. After pipeline
drain, the retired row's slot may be reused. The physical ring contains
$18$ memories of $512\times64$ bits, of which $480$ words cover an EO row.
Partially completed output strips are stored in a $2^{18}\times72$-bit
UltraRAM RowWindow. Packing four 16-bit EO pixels per word gives $256$
output-row slots; packing eight 8-bit IR samples gives $512$. The parameter
store is $32768\times144$ bits and the work queue $4096\times144$ bits.
These structures, ROMs, and metadata must all be included in memory claims;
the design is not a strict two-line architecture.

An output row becomes eligible only after all its required source rows
have retired. The emitter then blends seams and folds the panorama into
two HD raster bands. EO divides $3840$ columns into two $1920$-pixel halves.
IR divides its $3576$ valid columns into two $1788$-pixel halves and pads
each to $1920$ pixels. Both use $960$ active output rows within a
$1920\times1080$ raster. RowWindow alias guards stall the walker when the
emitter can advance; an unresolvable window overflow is reported and the
affected output row is blackened. It is not silently treated as a valid
stabilized image.

\subsection{Capture Coverage, Bandwidth, and Validation Boundary}
Each camera has four DDR source-bank address spaces. EO panorama capture
uses rows $0$--$1079$ and 16-pixel beats $8$--$111$, namely columns
$128$--$1791$. This fixed envelope includes every permitted transformed
segment endpoint, the sampler guard, and replay's eight-beat rounding for
all $2001$ roll values and pitch translations between $-98$ and $98$ rows.
Pitch monotonicity at fixed roll and affine variation within a segment
reduce the proof to endpoint tests at the two pitch extremes. IR retains
full $640\times512$ capture. A selected EO single-camera view requests full
capture and waits for a descriptor marked as adopting the full window;
an old narrow descriptor cannot satisfy it. The marker denotes the adopted
window, not an independent proof of a complete received raster.

Capture and output writes share DDR service. At an \emph{assumed} $30$
new frames/s, EO panorama capture requires $20.2176$ million useful
256-bit write beats/s, and output construction adds $3.456$, totaling
$23.6736$. IR panorama totals $1.8432+3.456=5.2992$ million beats/s.
The families' DDR captures are mutually gated in steady state. A historical
aggregate write measurement of $22.3$ million beats/s is below the EO
requirement and is not a measurement of this candidate. These write totals
exclude read service and do not establish sustainable frame rate. Output
scanout may repeat a completed image, so raster rate must be distinguished
from unique-content rate. A per-row capture ROM alone would not remove
this burden while replay requests its present global rectangular bounds.

Verification of the arithmetic, frame association, and renderer is reported in
Section~\ref{sec:verification}; measured and routed results follow in
Section~\ref{sec:hardware}.
